\section{Systemtests und Dauerbetrieb} \todo{7.4 + Abkürzungscheck}
Nach erfolgreicher Integration aller Komponenten erfolgt die Prüfung des Gesamtsystems im Dauerbetrieb.  
Dabei werden Stabilität, Temperaturverhalten und Zuverlässigkeit untersucht.  

Das System wird zunächst über einen Zeitraum von mehreren Stunden betrieben.  
Währenddessen erfolgt wiederholt die Kontrolle der Zeit-, Wetter- und Temperaturanzeigen.  
Die angezeigte Uhrzeit wird regelmäßig mit einer externen Referenz verglichen.  
Dabei kann keine sichtbare zeitliche Abweichung festgestellt werden.  
Zusätzlich wird überprüft, ob sich elektronische Komponenten oder Leitungen unzulässig erwärmen.  
Während des Betriebs treten keine kritischen Temperaturentwicklungen auf.  

Anschließend erfolgt ein Dauerbetrieb über drei Tage.  
Dabei wird die Funktion der Wortuhr in regelmäßigen Abständen kontrolliert.  
Alle Anzeigen arbeiten während dieses Zeitraums stabil und fehlerfrei.  
Nach erfolgreichem Abschluss der Dauerlauftests erfolgt der Aufbau der finalen Hardwareversion.  
Die zuvor verwendeten Breadboard- und Jumperleitungen werden durch fest verlötete Leitungen ersetzt.  
Zusätzlich werden die entwickelten und gedruckten Halterungen und Befestigungselemente montiert.  

Nach der Endmontage erfolgt eine erneute Überprüfung aller Funktionen und elektrischen Verbindungen.  
Dabei wird kontrolliert, ob alle Anzeigen korrekt arbeiten und sämtliche Leitungen ordnungsgemäß verbunden sind.  
Nach erfolgreichem Abschluss der Systemtests wird die Wortuhr an ihrem vorgesehenen Einsatzort montiert und in Betrieb genommen.